\section{Introduction}
\label{sec:intro}

Sign languages are natural languages with their own lexicon and grammar, used by millions of Deaf
and hard-of-hearing people. Automatic sign language processing covers tasks of increasing
difficulty: \emph{isolated} sign recognition (one sign per clip, closed vocabulary), continuous
recognition, and translation to and from spoken languages
\citep{bragg2019interdisciplinary,koller2020survey,yin2021including}. This report addresses
isolated recognition only. The distinction is not merely terminological: a classifier over a few
hundred glosses ignores the grammar of the language, its use of space and its non-manual markers,
and therefore cannot translate \citep{yin2021including}.

\emph{Pose-based} approaches first extract body, hand and face landmarks with an off-the-shelf
estimator such as MediaPipe Holistic
\citep{lugaresi2019mediapipe,bazarevsky2020blazepose,zhang2020mediapipehands,kartynnik2019facemesh},
then classify the resulting keypoint sequence. Compared with models operating on RGB video, they
are compact, largely insensitive to background and appearance, and allow inference without
transmitting video, which protects the privacy of users. The Google Isolated Sign Language
Recognition (ISLR) competition \citep{kaggle_islr} established this setting with
\numAslSamples{} landmark sequences of \numAslClasses{} American Sign Language (ASL) signs.

